%! AP Statistics — Unit 1, Lesson 1.9 — Guided Notes
\documentclass[10pt]{article}
\usepackage{apstats-article}
\usepackage{apstats-boxes}

%=============================================================================
\begin{document}

\pageheader{Unit 1, Lesson 1.9}{Guided Notes}
\namedateperiod

\vspace{0.10in}

% ── OBJECTIVE ────────────────────────────────────────────────────────────────
\begin{objectivebox}
\small By the end of this lesson, I will be able to\dots\\[4pt]
\writeline\\[1pt]
\writeline\\[1pt]
\writeline
\end{objectivebox}

\vspace{0.10in}

% ── VOCABULARY ───────────────────────────────────────────────────────────────
\begin{vocabbox}
\begin{multicols}{2}
\termblanklong{Parallel (side-by-side) boxplots}
\termblanklong{Back-to-back stemplot}
\termblanklong{Comparative statement}
\columnbreak
\termblanklong{Overlapping distributions}
\termblanklong{Explicit comparative language}
\termblanklong{Context (in a comparison)}
\end{multicols}
\end{vocabbox}

\vspace{0.10in}

% ── HOOK ─────────────────────────────────────────────────────────────────────
\begin{hookbox}
\small\textbf{Same data, same scale --- now we can compare}

\vspace{5pt}
The parallel boxplots below show the same ACT score data from the warm-up.
Note how placing both on one number line makes comparisons immediate.

\vspace{4pt}
\begin{center}
\begin{tikzpicture}[scale=1.0]
  \draw[thick] (0.0, 0) -- (6.80, 0);
  \foreach \v/\x in {14/0.00, 16/0.60, 18/1.20, 20/1.80, 22/2.40,
                      24/3.00, 26/3.60, 28/4.20, 30/4.80, 32/5.40, 34/6.00, 36/6.60} {
    \draw (\x, -0.12) -- (\x, 0.12);
    \node[below, font=\small] at (\x, -0.12) {\v};
  }
  \node[below, font=\small\itshape] at (3.30, -0.56) {ACT composite score};
  \node[left, font=\small] at (0.0, 1.10) {No Prep};
  \draw[thick, navy] (0.30, 1.10) -- (1.50, 1.10);
  \draw[thick, navy] (1.50, 0.85) rectangle (3.30, 1.35);
  \draw[thick, navy] (2.40, 0.85) -- (2.40, 1.35);
  \draw[thick, navy] (3.30, 1.10) -- (4.50, 1.10);
  \draw[thick, navy] (0.30, 0.95) -- (0.30, 1.25);
  \draw[thick, navy] (4.50, 0.95) -- (4.50, 1.25);
  \node[left, font=\small] at (0.0, 0.35) {Test Prep};
  \draw[thick, navy] (1.80, 0.35) -- (3.00, 0.35);
  \draw[thick, navy] (3.00, 0.10) rectangle (4.80, 0.60);
  \draw[thick, navy] (3.90, 0.10) -- (3.90, 0.60);
  \draw[thick, navy] (4.80, 0.35) -- (6.30, 0.35);
  \draw[thick, navy] (1.80, 0.20) -- (1.80, 0.50);
  \draw[thick, navy] (6.30, 0.20) -- (6.30, 0.50);
\end{tikzpicture}
\end{center}

\vspace{4pt}
What conclusions can you see at a glance? What still requires calculation?
\end{hookbox}

\vspace{0.08in}

% ── SECTION 1: READING PARALLEL BOXPLOTS ─────────────────────────────────────
\begin{notesbox}{Reading Parallel Boxplots}
\small
\begin{multicols}{2}

\textbf{\textcolor{navy}{What to read from each boxplot:}}
\begin{itemize}[itemsep=3pt]
  \item Median $=$ \blank{3cm} (center)
  \item $Q_1, Q_3$ and IQR $=$ \blank{1.5cm} (spread)
  \item Min, max, range $=$ \blank{1.5cm}
  \item Outliers (individual \blank{1.5cm})
  \item Approximate shape (\blank{2.5cm})
\end{itemize}

\vspace{6pt}
\textbf{\textcolor{navy}{From the ACT hook:}}

\vspace{4pt}
\renewcommand{\arraystretch}{1.4}
\begin{tabularx}{\columnwidth}{@{} l X X @{}}
& No Prep & Test Prep \\
\midrule
Median & \blank{1.4cm} & \blank{1.4cm} \\
$Q_1$  & \blank{1.4cm} & \blank{1.4cm} \\
$Q_3$  & \blank{1.4cm} & \blank{1.4cm} \\
IQR    & \blank{1.4cm} & \blank{1.4cm} \\
Range  & \blank{1.4cm} & \blank{1.4cm} \\
\end{tabularx}

\columnbreak

\textbf{\textcolor{navy}{Comparing across groups:}}
\begin{itemize}[itemsep=4pt]
  \item \textit{Center}: Compare the two \blank{2cm}.
  \item \textit{Spread}: Compare the two \blank{2cm}
        (or ranges, or both).
  \item \textit{Overlap}: How much of each IQR box
        overlaps the other? Large overlap means the groups
        are \blank{2.5cm}; small overlap suggests
        a \blank{3cm} difference.
  \item \textit{Shape \& Outliers}: Note any skewness or
        outlier dots visible in each row.
\end{itemize}

\vspace{6pt}
\textit{Important:} The same \blank{3cm} must be
used for both boxplots. Separate scales make
comparison \blank{2cm}.

\end{multicols}
\end{notesbox}

\vspace{0.08in}

% ── SECTION 2: AP COMPARATIVE LANGUAGE ───────────────────────────────────────
\begin{notesbox}{Writing AP Comparative Statements}
\small
\begin{multicols}{2}

\textbf{\textcolor{navy}{The AP scoring rule:}}\\
A comparative statement earns credit only if it:
\begin{enumerate}[itemsep=3pt, label=\alph*.]
  \item Uses a \blank{2.5cm} word
        (higher, lower, more, less, similar, \ldots)
  \item Names \blank{2cm} groups with context
  \item States \blank{2cm} values
\end{enumerate}

\vspace{6pt}
\textbf{\textcolor{navy}{Non-example (no credit):}}\\
\textit{``The median ACT score for the Test Prep group
was 27.''}\\[3pt]
Problem: \blank{6cm}

\vspace{6pt}
\textbf{\textcolor{navy}{Example (full credit):}}\\
\textit{``The median ACT score for the Test Prep group
(27) was \textbf{5 points higher than} the median for the
No Prep group (22).''}

\columnbreak

\textbf{\textcolor{navy}{The comparison sentence frame:}}

\vspace{4pt}
\begin{tcolorbox}[colback=goldbg, colframe=goldacc, boxrule=0.6pt,
    arc=2mm, left=3mm, right=3mm, top=2mm, bottom=2mm]
\small
``The \blank{2cm} [variable] for
[Group A] (\blank{1.5cm}) was
\blank{2cm} [higher/lower] than
the \blank{2cm} for [Group B]
(\blank{1.5cm}).''
\end{tcolorbox}

\vspace{6pt}
\textbf{\textcolor{navy}{Full AP comparison uses SOCS:}}
\begin{itemize}[itemsep=3pt]
  \item \textbf{S}hape: Describe shape and skewness for each.
  \item \textbf{O}utliers: Note any outlier dots; verify if asked.
  \item \textbf{C}enter: Compare medians (required).
  \item \textbf{S}pread: Compare IQRs or ranges (required).
\end{itemize}

\end{multicols}
\end{notesbox}

\vspace{0.08in}

% ── SECTION 3: BACK-TO-BACK STEMPLOTS ────────────────────────────────────────
\begin{notesbox}{Back-to-Back Stemplots}
\small
\begin{multicols}{2}

\textbf{\textcolor{navy}{When to use:}}\\
Best for \blank{2cm} datasets ($n < 30$) when you
want to see the shape of \blank{2cm} groups.

\vspace{6pt}
\textbf{\textcolor{navy}{Structure:}}\\
A single \blank{2cm} column sits in the middle.
Leaves for Group 1 extend to the \blank{1.5cm};
leaves for Group 2 extend to the \blank{1.5cm}.
Group 1 leaves are written in \blank{2.5cm} order
(largest to smallest reading outward).

\vspace{6pt}
\textbf{\textcolor{navy}{Reading centers and spread:}}
\begin{itemize}[itemsep=3pt]
  \item Median: count to the \blank{2.5cm} leaf.
  \item Spread: read min and max from the ends.
  \item Shape: look at where leaves are
        \blank{2cm} or \blank{2cm}.
\end{itemize}

\columnbreak

\textbf{\textcolor{navy}{ACT example --- back-to-back stemplot}}\\
\textit{(Key: $1 \mid 8 = 18$ points; leaf unit = 1)}

\vspace{4pt}
\begin{center}
\small
\renewcommand{\arraystretch}{1.3}
\begin{tabular}{r c l}
\multicolumn{1}{c}{\textit{No Prep}} & Stem & \textit{Test Prep} \\
\hline
\hspace{0.5em}8\;7 & 1 & \\
8\;7\;5\;3\;2\;0 & 2 & 2\;4\;5\;6\;7\;8\;9 \\
 & 3 & 0\;3 \\
\hline
\end{tabular}
\end{center}

\vspace{6pt}
\small No Prep ($n=8$): 17, 18, 20, 22, 23, 25, 26, 27\\
Test Prep ($n=9$): 22, 24, 25, 26, 27, 28, 29, 30, 33

\vspace{4pt}
From the stemplot:
\begin{itemize}[itemsep=2pt]
  \item No Prep median $= \blank{2cm}$ (4th value)
  \item Test Prep median $= \blank{2cm}$ (5th value)
  \item Which group has more spread? \blank{3cm}
\end{itemize}

\end{multicols}
\end{notesbox}

\vspace{0.08in}

% ── QUICK CHECK ──────────────────────────────────────────────────────────────
\begin{tcolorbox}[colback=goldbg, colframe=goldacc, boxrule=0.8pt, arc=2mm,
    left=4mm, right=4mm, top=2mm, bottom=2mm]
\small\textbf{Quick Check:} Using the parallel boxplots in the hook, write one
AP-quality comparative statement about \textbf{spread} for the two ACT groups.
Your sentence must name both groups, use a comparison word, give both IQR
values, and include context (ACT scores).\\[8pt]
\writeline\\[1pt]\writeline
\end{tcolorbox}

\end{document}
